\documentclass[12pt]{article}
\usepackage[utf8]{inputenc}
\usepackage[T1]{fontenc}
\usepackage[margin=1in]{geometry}
\usepackage{setspace}
\usepackage{xcolor}
\usepackage[colorlinks=true, urlcolor=blue!60!black, linkcolor=blue!60!black]{hyperref}
\setlength{\parindent}{2em}
\setlength{\parskip}{1.25em}
\renewcommand{\baselinestretch}{1.0}
\begin{document}
\begin{center}
{\LARGE \textit{Stuart v. Laird}, 5 U.S. (1 Cranch) 299 (1803)\par}
\vspace{0.8em}
{\large Case Brief --- Four-Page Edition\par}
\end{center}
\vspace{1em}

\section{Memory Jogger}

Six days after \textit{Marbury v. Madison}, a unanimous Court avoided a showdown with the Jeffersonian Republicans, holding that Congress may transfer pending cases between reorganized lower federal courts and that longstanding practice settled the constitutionality of Supreme Court Justices riding circuit without separate commissions (\href{https://supreme.justia.com/cases/federal/us/5/299/}{Justia}).

\section{Detailed Case Facts}

In January 1801 John Laird, for the Scottish firm Laird \& Robertson, sued Hugh Stuart for an unpaid debt in the federal circuit court for Virginia created by the Judiciary Act of 1789. In February 1801 the lame-duck Federalist Congress passed the Judiciary Act of 1801, the ``Midnight Judges Act,'' which abolished circuit riding, created sixteen life-tenured circuit judgeships, and transferred pending cases to the new courts. Laird won judgment in the new Fourth Circuit court, and Stuart and his surety, Charles L. Carter, gave a forthcoming bond for levied property.

The new Democratic-Republican Congress then passed the Repeal Act of March 8, 1802, abolishing the sixteen judgeships, and the Judiciary Act of April 29, 1802, which reinstated circuit riding by Supreme Court Justices, transferred pending proceedings to reconstituted circuit courts, and canceled the Court's June 1802 Term. In the reconstituted Fifth Circuit, Chief Justice Marshall, sitting as circuit justice, entered judgment for Laird on the bond. Represented by Charles Lee, who also represented William Marbury, Stuart challenged both the transfer and the authority of Justices to sit as circuit judges without distinct commissions.

\section{Procedural History}

The suit began in the Middle Circuit (1801), was transferred to the Fourth Circuit under the 1801 Act, and then to the Fifth Circuit under the 1802 Act, where Marshall entered judgment on the bond. On writ of error, the Supreme Court heard argument on February 23--24, 1803, and affirmed on March 2, 1803.

\section{Judicial Votes}

Affirmed, \textbf{5--0}.
\begin{itemize}
\item \textbf{Majority:} Justice Paterson, joined by Justices Cushing, Chase, Washington, and Moore.
\item \textbf{Concurrences and dissents:} None.
\item \textbf{Not participating:} Chief Justice Marshall, who had tried the cause below.
\end{itemize}

\section{Holding}

Under Article III, Congress may establish inferior courts from time to time and transfer pending cases, including ancillary proceedings such as forthcoming bonds, from one inferior tribunal to another. Practice and acquiescence since the founding of the judicial system establish that Justices may sit as circuit judges under their Supreme Court commissions. The opinion never addressed Lee's central argument that the Repeal Act unconstitutionally stripped the sixteen judges of their offices and salaries under the Good Behavior Clause, though the decision has long been treated as implicitly validating Congress's power to abolish lower courts.

\section{Analysis of Opinions}

\textbf{Paterson (for the Court).} In barely two pages, Paterson disposed of both assignments of error. First, absent statute only the issuing court could enforce the bond, but Congress's power to create inferior tribunals includes power to transfer causes between them, and ``there are no words in the constitution to prohibit or restrain the exercise of legislative power.'' The bond was ``an appendage to the cause, or rather a component part of the proceedings,'' and the 1802 Act authorized the Fifth Circuit to carry it into effect.

Second, on the objection that Justices could not sit on circuit without separate commissions under the Appointments Clause and Article III's allocation of jurisdiction, Paterson bypassed textual analysis and invoked early institutional practice---what modern scholars call constitutional liquidation or historical gloss. Circuit riding since 1789 was a ``contemporary interpretation of the most forcible nature'' that put the question ``at rest.''

\textbf{No separate opinions.} Marshall and Chase had privately doubted that Justices could ride circuit under their existing commissions, but the Justices closed ranks behind Paterson to preserve the Court amid threats of Jeffersonian impeachments.

\section{Examples: Future Applications}

\begin{enumerate}
\item \textbf{Same-Side (Valid): Splitting the Ninth Circuit.} Congress creates a new Twelfth Circuit and transfers pending Arizona and Nevada appeals to it; reorganization and docket transfer are squarely permitted.
\item \textbf{Same-Side (Valid): Sitting by Designation.} District judges and retired Justices sit on courts of appeals without new nominations, resting on statute and generations of practice, like circuit riding.
\item \textbf{Opposite-Side (Invalid): Abolishing One Judge's Seat.} Abolishing a single district seat to oust an unpopular judge and end her salary, while reassigning her docket to an identical new court, violates the Good Behavior and Compensation Clauses; \textit{Stuart} did not decide an ousted judge's claim.
\item \textbf{Opposite-Side (Invalid): Transfer to an Executive Tribunal.} Moving pending private contract suits to an agency tribunal whose officers serve at the President's pleasure takes cases out of Article III altogether.
\item \textbf{Fence-Sitter: Eighteen-Year Term Limits.} A law moving Justices to ``senior Justice'' status after 18 years: proponents invoke Congress's power over assignments; opponents say it removes Justices from the ``one supreme Court'' without impeachment and lacks founding-era acquiescence.
\end{enumerate}

\section{Critique}

\textbf{Originalist views.} William Baude and other students of liquidation treat \textit{Stuart} as foundational proof that early practice can settle indeterminate text, a principle later cited in \textit{NLRB v. Noel Canning}. Strict textualists object that the Court used barely thirteen years of reluctant, protested circuit riding to override serious textual objections, and left unresolved whether Congress can evade life tenure by abolishing a judge's court.

\textbf{Pragmatic views.} Bruce Ackerman and historical-gloss scholars such as Curtis Bradley and Neil Siegel read \textit{Stuart} with \textit{Marbury} as judicial statecraft: the Court asserted judicial review while yielding to the Republicans' core statutory priority, and sensibly left Congress room to adapt the courts.

\section{Key Quotations}

\begin{quote}
``To this objection, which is of recent date, it is sufficient to observe, that practice and acquiescence under it for a period of several years, commencing with the organization of the judicial system, affords an irresistible answer, and has indeed fixed the construction.''
\end{quote}
\noindent --- Justice Paterson, for the Court, 5 U.S. (1 Cranch) at 309.

\begin{quote}
``It is a contemporary interpretation of the most forcible nature. This practical exposition is too strong and obstinate to be shaken or controlled. Of course, the question is at rest, and ought not now to be disturbed.''
\end{quote}
\noindent --- Justice Paterson, for the Court, 5 U.S. (1 Cranch) at 309.

\begin{quote}
``Congress have constitutional authority to establish from time to time such inferior tribunals as they may think proper, and to transfer a cause from one such tribunal to another. In this last particular, there are no words in the constitution to prohibit or restrain the exercise of legislative power.''
\end{quote}
\noindent --- Justice Paterson, for the Court, 5 U.S. (1 Cranch) at 309.

\end{document}
